\begin{theorem}[Monogenic uniformizers]
\label{mod:localfield-monogenicuniformizer}
Let $K/F$ be a finite valuative extension of nonarchimedean local fields and
assume
\[
  f(K/F)=\operatorname{residueDegree}(F,K)=1.
\]
For every integral uniformizer
$\pi_K\in\mathcal O_K$, prove
\[
 \operatorname{Algebra.adjoin}_{\mathcal O_F}\{\pi_K\}
   =\mathcal O_K.
\]
There is consequently an integral uniformizer with its uniformizer proof and
this monogenicity equality:
\[
 \exists\,\pi_K\in\mathcal O_K,\quad
 \operatorname{IsUniformizer}(\pi_K)\ \wedge\
 \operatorname{Algebra.adjoin}_{\mathcal O_F}\{\pi_K\}
   =\mathcal O_K.
\]

The result is a reusable local-field theorem.  Its statement and proof contain
no cyclicity, Galois, ramification-break, character, local-constant, or case
hypotheses.
\LeanFile{LanglandsFirstMainLemma/LocalField/MonogenicUniformizer.lean}
\lean{LanglandsFirstMainLemma.algebra_adjoin_uniformizer_eq_top_of_residueDegree_eq_one}
\lean{LanglandsFirstMainLemma.monogenicUniformizer}
\leanok
\SourceText{The totally ramified uniformizer input used in the proof of Theorem thm:norm-filtration and in Section sec:first-main-completion.}
\uses{mod:localfield-valuation,mod:localfield-lattices,mod:localfield-residuefield,mod:localfield-extension}
\FormalizationContract
Put $C=\mathcal O_F[\pi_K]$.  Residue degree one makes the canonical map
$k_F\to k_K$ surjective.  Hence every $b\in\mathcal O_K$ has a lower integral
lift modulo the maximal ideal, and principality
$\mathfrak p_K=(\pi_K)$ gives
\[
 \mathcal O_K=C+\pi_K\mathcal O_K.
\]
The finite $\mathcal O_F$-module $\mathcal O_K$ is also finite over $C$.
Furthermore, integrality and lying-over show that every maximal ideal of $C$
lifts to the unique maximal ideal of $\mathcal O_K$; consequently
$\pi_K\in\operatorname{Jac}(C)$.  Nakayama's lemma applied over $C$ proves
$C=\mathcal O_K$.  This finite argument replaces an infinite uniformizer
expansion and uses neither a prior monogenicity assertion nor a power basis.

The existential statement is obtained from the universal theorem and the
existing canonical-valuation uniformizer-existence theorem.  In particular,
the proof does not choose a specially adapted generating uniformizer.
\end{theorem}
